% TOP_PROBLEM: {"id": 4700007, "title": "Trigonometric Abel differential equations II", "category_id": 11, "category": {"id": 11, "name": "dynamical_systems", "display_name": "Dynamical Systems", "description": "Problems about long-term behavior of deterministic systems, Hamiltonian dynamics, and periodic orbits.", "slug": "dynamical-systems", "order_index": 11, "created_at": "2026-07-31T15:26:25.671Z"}, "status": "open", "problem_number": "AMR-046-0007", "set_id": 13}
% FIRST_POSTED: 2026-10-02
% ENTRY_KIND: statement_only
% UnsolvedMath ID 4700007; source catalogue code AMR-046-0007
% !TeX program = lualatex
% Definition and sources only; this file is not an LLM solution attempt.
\documentclass[11pt,a4paper]{article}
\usepackage[margin=25mm]{geometry}
\usepackage{fontspec}
\setmainfont{lmroman10-regular.otf}[BoldFont=lmroman10-bold.otf,
 ItalicFont=lmroman10-italic.otf,BoldItalicFont=lmroman10-bolditalic.otf]
\usepackage{amsmath,amssymb,mathtools}
\usepackage{unicode-math}
\setmathfont{latinmodern-math.otf}
\AtBeginDocument{\let\setminus\smallsetminus}
\usepackage{xurl,hyperref}
\hypersetup{colorlinks=true,urlcolor=blue}
\setcounter{secnumdepth}{0}
\setlength{\parindent}{0pt}
\setlength{\parskip}{5pt}
\setlength{\emergencystretch}{3em}
\begin{document}
\section*{Trigonometric Abel differential equations II}\label{4700007}
{\small UnsolvedMath ID 4700007 · AMR-046-0007 · Dynamical Systems · AMR Open Problem Lists}
\subsection{Definitions and mathematical statement}
Fix integers $p>q\ge2$ and nonnegative integers $m,n$. A real trigonometric
polynomial of degree $m$ is a function
$A(t)=a_0+\sum_{j=1}^m(a_j\cos jt+\widetilde a_j\sin jt)$,
with $(a_m,\widetilde a_m)\ne(0,0)$ when $m>0$ and $a_0\ne0$ when $m=0$.
Define $B$ similarly with degree $n$, and consider
\[
 x'=A(t)x^p+B(t)x^q.
\]
The degrees are exactly $m$ and $n$, as in the source question; the leading
harmonics therefore remain present throughout this fixed-degree family.

A periodic solution here is a real-valued solution that exists without
blow-up on all of $\mathbb R$ and satisfies $x(t+2\pi)=x(t)$. A limit cycle
means that its value at $t=0$ is an isolated fixed point of the period return
map. The question asks for the maximum number of these cycles as a function
of $p,q,m,n$, with all coefficients otherwise unrestricted.

The constant zero solution is included if isolated. Nonzero solutions can
be positive or negative and both signs must be counted; their parity
behavior depends on the exponents. Distinct initial values at the fixed
forcing phase represent distinct periodic solutions. A period which divides
$2\pi$ is allowed. Continuous families of periodic solutions and solutions
that become infinite in finite time do not contribute isolated cycles.
The target is a uniform sharp bound and examples attaining it, or a proof
that no finite bound exists for some fixed degree and exponent data.
\subsection{Short English statement}
Determine the maximum number of periodic limit cycles from the two exponents and the trigonometric coefficient degrees.
\subsection{Sources}
\begin{itemize}
\item[S1] ULAM.AI and the UnsolvedMath Contributors, supplied problems.json, record 4700007 (AMR-046-0007), AMR Open Problem Lists. Catalogue identity and source transcription; CC BY 4.0.
\url{https://huggingface.co/datasets/ulamai/UnsolvedMath}.
\item[S2] Gasull - Some open problems in low dimensional dynamical systems (2020), Problem environment 7: Trigonometric Abel differential equations II. Original source indexed by AMR-046-0007.
\url{https://arxiv.org/abs/2012.02524}.
\end{itemize}
\end{document}
